
%\needspace{4\baselineskip}
\color{uniblue}{\subsection{Защита от феррорезонанса (ЗФР)}}
\color{black}

\begin{enumerate}[label=\arabic{section}.\arabic{subsection}.\arabic*, labelsep=4pt, leftmargin=0pt, itemindent=57pt]

\item
Защита от феррорезонанса предназначена для защиты обмоток ТН от перенапряжений и токовых перегрузок в режиме феррорезонанса. В устройстве защита от феррорезонанса представлена функцией «ЗФР».

\item
Функция фиксирует наличие режима феррорезонанса по повышению напряжения 3U0 выше заданной уставки и, через выдержку времени, определяемую уставкой таймера Т1, выдает управляющее воздействие на выходное реле управления режимом специальной обмоткой ТН (дешунтируя или шунтируя ее контактом выходного реле).

\item
При повышении напряжения НП выше уставки и отсутствии блокирующих сигналов происходит пуск защиты. По истечении выдержки времени Т1, в течение которой величина напряжения НП постоянно была выше уставки срабатывания, происходит срабатывание защиты. Сброс сигнала срабатывания функции осуществляется при следующих условиях:
\begin{itemize}
    \item	при снижении величины напряжения $3U_0$ ниже уставки возврата ИО;
    \item	при появлении сигнала сброса от кнопки «Сброс»;
    \item	при появлении сигнала неисправности цепей напряжения БНН при введенной уставке SGF2.
\end{itemize}

\item
Ввод функции в работу на этапе параметрирования устройства осуществляется программным переключателем SGF1 «Ввод\_функции». Информация о введенной функции отображается выходным сигналом «ЗФР: Ввод».

\item\sloppy
Функция может быть выведена из работы оперативно путем активации сигнала «ОВ ЗФР». Выведенное состояние функции характеризуется активным состоянием сигнала «ЗФР: Оперативный вывод».

\item
При активном состоянии входного сигнала «ЗФР на сигн» (ключ «ЗФР на сигн» введен) действие функции переводится на сигнал.

\item
Алгоритм работы функции выполнен в соответствии с рисунком \ref{zfr:img1}. В таблице \ref{zfr:tbl1} приведены параметры, необходимые для настройки функции.

\vspace{3mm}
\begin{figure}[H]
\centering
\includegraphics[width=1\textwidth,height=1\textheight,keepaspectratio]{\fbpath/202. ЗФР/_latex/img/zfr.pdf}
\captionof{figure}{Функциональная схема логики ЗФР}\label{zfr:img1}
\end{figure}

\small
\begin{longtable}{|>{\centering\arraybackslash}m{5.3cm}|>{\centering\arraybackslash}m{3.3cm}|>{\centering\arraybackslash}m{4.2cm}|>{\centering\arraybackslash}m{1.8cm}|>{\centering\arraybackslash}m{1cm}|}
\caption{Параметры для настройки ЗФР\hfill\vspace{-0.5\baselineskip}}\label{zfr:tbl1}\\ 
\hline
\rowcolor{gray!30}
Параметр (Параметр на ИЧМ) & Условное обозначение на схеме & Значение/ Диапазон & Единица измерения & Шаг \\ 
\hline
\endfirsthead
\caption*{\hspace{3pt}\emph{Продолжение таблицы \ref{zfr:tbl1}\hfill\vspace{-0.5\baselineskip}}} \\ % сделано по ГОСТ 2.105 п.6.8.7
\hline
\rowcolor{gray!30}
Параметр (Параметр на ИЧМ) & Условное обозначение на схеме & Значение/ Диапазон & Единица измерения & Шаг \\ 
%\hline
\endhead
%\hline
\endfoot
\hline
\endlastfoot
%==+t1*RESOPTOV1|VT_RESOTOV
\centering Ввод функции в работу (Ввод\_функции) & \centering SGF1 & \centering 0 = Не предусмотрено\\1 = Предусмотрено & \centering -- & \centering \arraybackslash -- \\
\hline
\centering Напряжение срабатывания НП (3U0ср) & \centering 3U0> & \centering 2,0 ... 150,0 & \centering В & \centering \arraybackslash 0,1 \\
\hline
\centering Напряжение возврата НП (3U0в) & \centering 3U0< & \centering 2,0 ... 100,0 & \centering В & \centering \arraybackslash 0,1 \\
\hline
\centering Режим контроля от БНН (Реж\_БНН) & \centering SGF2 & \centering 0 = Без контроля\\1 = С контролем & \centering -- & \centering \arraybackslash -- \\
\hline
\centering Выдержка времени срабатывания (Tср) & \centering T1 & \centering 0,00 ... 100,00 & \centering с & \centering \arraybackslash 0,01 \\
%===t1
\end{longtable}
\normalsize

\end{enumerate}
